% Reader question: when a conflict needs a decision above the people involved,
% when is there exactly one place to send it? Claim: escalation has a next hop
% only when the principals above the conflicting object form one line; a
% diamond (two parents for one scope) leaves no next hop. Evidence: the
% chapter's unity-of-command theorem and its exercise: Bob holds the charter
% and delegates to Alice and to tests/; in the diamond he also delegates to
% Carol, and Alice and Carol each delegate payments/. Must distinguish: the
% one-line case (a forced route upward) from the diamond (two incomparable
% parents); principals from objects. Grammar: a before/after pair of
% delegation graphs with identical node positions, the right panel adding
% Carol and her edge; rejected: one graph annotated with both cases (the
% reader has to subtract) and an org chart without objects (the conflict sits
% on an object, not on a person). Five-second test: "left, the conflict climbs
% Alice then Bob; right, it could go to Alice or Carol, so it has nowhere to
% go."
\begin{figure}[H]
\centering
\pdfigurehue{pdteal}%
\begin{tikzpicture}[pd figure,x=1cm,y=1cm,
  who/.style={pd state,minimum width=1.15cm},
  object/.style={pd artifact,minimum width=1.2cm,font=\pdfiglabelmono},
  owns/.style={pd arrow,shorten <=2pt,shorten >=2pt},
  goesup/.style={pd focus arrow,shorten <=3pt,shorten >=3pt},
  nohop/.style={pd breach arrow,shorten <=3pt,shorten >=3pt},
  query/.style={pd breach label,font=\pdfiglabelbold}]
  % Both panels share one grid: Bob (2.6,3.5); Alice (.8,1.85), Carol
  % (2.6,1.85), tests/ (4.55,1.85); payments/ (1.7,0). Only Carol and her edges
  % are new on the right. The escalation route runs just inside the
  % delegation edges it climbs, so it never crosses them.
  % ---- left: one line of authority ------------------------------------------
  \node[pd title,anchor=west] at (0,4.30) {One line:\ escalate to Bob};
  \node[who] (bobL) at (2.6,3.5) {Bob};
  \node[who] (aliceL) at (.8,1.85) {Alice};
  \node[object] (testsL) at (4.55,1.85) {tests/};
  \node[object] (payL) at (1.7,0) {payments/};
  \draw[owns] (bobL) -- (aliceL);
  \draw[owns] (bobL) -- (testsL);
  \draw[owns] (aliceL) -- (payL);
  \draw[goesup] (payL.north) to[bend right=14] (aliceL.south east);
  \draw[goesup] (aliceL.north east) to[bend right=12] (bobL.south);
  \node[pd focus label,anchor=west] at (2.05,1.0) {conflict goes up};
  % ---- divider ---------------------------------------------------------------
  \draw[pd hairline] (5.45,-.40) -- (5.45,4.50);
  % ---- right: a diamond, no next hop ---------------------------------------
  \begin{scope}[xshift=5.75cm]
    \node[pd title,anchor=west] at (0,4.30) {Two lines:\ no next hop};
    \node[who] (bobR) at (2.6,3.5) {Bob};
    \node[who] (aliceR) at (.8,1.85) {Alice};
    \node[pd focus outline,minimum width=1.15cm] (carolR) at (2.6,1.85) {Carol};
    \node[object] (testsR) at (4.55,1.85) {tests/};
    \node[object] (payR) at (1.7,0) {payments/};
    \draw[owns] (bobR) -- (aliceR);
    \draw[owns] (bobR) -- (carolR);
    \draw[owns] (bobR) -- (testsR);
    \draw[owns] (aliceR) -- (payR);
    \draw[owns] (carolR) -- (payR);
    \draw[nohop] (payR.north) -- (aliceR.south east);
    \draw[nohop] (payR.north) -- (carolR.south west);
    \node[query] at (1.7,1.85) {?};
    \node[pd label,anchor=north,text width=4.6cm] at (1.7,-.45)
      {Alice and Carol both delegated \texttt{payments/}};
  \end{scope}
\end{tikzpicture}
\caption{Escalation has a next hop only when the people above the conflict
form one line. Fix the diamond by giving \texttt{payments/} one parent,
splitting it, or fixing a tie-break order. [constructed example]}
\label{fig:audit-escalate}
\end{figure}
